% Apéndice: Kernel Sources
% Codificación: UTF-8

\chapter{Kernel Sources}
\label{app:kernels}

This appendix lists the OpenCL C source of every kernel of the suite of
Section~\ref{sec:eval-results}, in the order of Table~\ref{tab:suite}.
The shuffle kernels carry the \texttt{intel\_reqd\_sub\_group\_size}
attribute, which fixes the sub-group size at 32, and take the number of
elements per work-item as the argument \texttt{N}.

\section*{Basics}

\begin{lstlisting}[language=C, caption={[Source of vadd]OpenCL source of \texttt{vadd}.}, label={lst:app-vadd}]
__kernel void vadd(__global int *a, __global int *b, __global int *c)
{
    int i = get_global_id(0);
    c[i] = a[i] + b[i];
}
\end{lstlisting}

\section*{Control flow}

\begin{lstlisting}[language=C, caption={[Source of vadd\_if\_else]OpenCL source of \texttt{vadd\_if\_else}.}, label={lst:app-vadd-if-else}]
__kernel void vadd_if_else(__global int *a, __global int *b, __global int *c)
{
    int i = get_global_id(0);
    if (i % 4 == 0)
        c[i] = a[i] + b[i];
    else
        c[i] = a[i] - 1;
}
\end{lstlisting}

\begin{lstlisting}[language=C, caption={[Source of vadd\_mod]OpenCL source of \texttt{vadd\_mod}.}, label={lst:app-vadd-mod}]
__kernel void vadd_mod(__global int *a, __global int *b, __global int *c)
{
    int i = get_global_id(0);
    if (i % 2 == 0) {
        c[i + 1] = a[i + 1] + b[i];
    }
}
\end{lstlisting}

\begin{lstlisting}[language=C, caption={[Source of nested\_if]OpenCL source of \texttt{nested\_if}.}, label={lst:app-nested-if}]
__kernel void nested_if(__global int *a, __global int *b, __global int *c)
{
    int j = get_local_id(0);
    int i = a[j];
    int r = -1;
    if (i >= 16) {
        if (i >= 24) {
            if (i >= 28) {
                if (i >= 30)
                    r = 30;
                else
                    r = 28;
            }
            else if (i >= 26)
                r = 26;
            else
                r = 24;
        }
        else {
            if (i >= 20) {
                if (i >= 22)
                    r = 22;
                else
                    r = 20;
            }
            else if (i >= 18)
                r = 18;
            else
                r = 16;
        }
    }
    else {
        if (i >= 8) {
            if (i >= 12) {
                if (i >= 14)
                    r = 14;
                else
                    r = 12;
            }
            else if (i >= 10)
                r = 10;
            else
                r = 8;
        }
        else {
            if (i >= 4) {
                if (i >= 6)
                    r = 6;
                else
                    r = 4;
            }
            else if (i >= 2)
                r = 2;
            else
                r = 0;
        }
    }
    c[j] = r;
}
\end{lstlisting}

\begin{lstlisting}[language=C, caption={[Source of loop]OpenCL source of \texttt{loop}.}, label={lst:app-loop}]
__kernel void loop(__global const int *a, __global int *out)
{
    int gid  = get_global_id(0);
    int lid  = get_sub_group_local_id();
    int base = gid - lid;
    int s = 0;
    for (int i = 0; i < lid; i++)
        s += a[base + i];
    out[gid] = s;
}
\end{lstlisting}

\begin{lstlisting}[language=C, caption={[Source of nested\_loop]OpenCL source of \texttt{nested\_loop}.}, label={lst:app-nested-loop}]
__kernel void nested_loop(__global const int *a, __global int *out)
{
    int gid  = get_global_id(0);
    int lid  = get_sub_group_local_id();
    int base = gid - lid;
    int s = 0;
    __attribute__((opencl_unroll_hint(1)))
    for (int i = 0; i < lid; i++) {
        __attribute__((opencl_unroll_hint(1)))
        for (int j = 0; j < 4; j++)
            s += a[base + 4 * i + j];
    }
    out[gid] = s;
}
\end{lstlisting}

\begin{lstlisting}[language=C, caption={[Source of control\_flow]OpenCL source of \texttt{control\_flow}.}, label={lst:app-control-flow}]
__kernel void control_flow(__global int *result)
{
    int lid = get_sub_group_local_id();
    int cnt = 0;
    for (int i = 0; i < lid; i++) {
        int k = i % 8;
        switch (k) {
        case 0:
            cnt++;
            break;
        case 1:
            cnt += i;
        case 2:
            cnt--;
            break;
        case 3:
        case 4:
            cnt += 2;
            continue;
            break;
        case 5:
            {
            int j = i;
            do {
                cnt += j;
            } while (j-- > 0);
            break;
            }
        case 7:
            {
            int m = i / 8;
            if (m == 0)
                cnt--;
            else if (m == 1)
                break;
            else if (m == 2)
                continue;
            else
                cnt++;
            }
        default:
            cnt++;
        }
    }
    result[lid] = cnt;
}
\end{lstlisting}

\begin{lstlisting}[language=C, caption={[Source of control\_flow1]OpenCL source of \texttt{control\_flow1}.}, label={lst:app-control-flow1}]
__kernel void control_flow1(__global int *result)
{
    int lid = get_sub_group_local_id();
    int cnt = 0;
    for (int i = 0; i < 2 * lid; i++) {
        int k = i % 8;
        switch (k) {
        case 0:
            switch (i) {
            case 0:
                break;
            case 8:
                cnt++;
            case 16:
                continue;
                break;
            case 24:
                cnt--;
                break;
            default:
                for (int j = 0; j < i; j++) {
                    if (j < 4) {
                        cnt++;
                    }
                    else if (j == 8)
                        cnt += 3;
                    else if (j > 16 && j < 24)
                        continue;
                    else
                        cnt += 2;
                }
            }
            break;
        case 1:
            cnt += i;
        case 2:
            cnt--;
            break;
        case 3:
        case 4:
            cnt += 2;
            continue;
            break;
        case 5:
            {
            int j = i;
            do {
                cnt += j;
            } while (j-- > 0);
            break;
            }
        case 7:
            {
            int m = i / 8;
            if (m == 0)
                cnt--;
            else if (m == 1)
                break;
            else if (m == 2)
                continue;
            else
                cnt++;
            }
        default:
            cnt++;
        }
    }
    result[lid] = cnt;
}
\end{lstlisting}

\section*{Sub-group communication}

\begin{lstlisting}[language=C, caption={[Source of shfl1]OpenCL source of \texttt{shfl1}.}, label={lst:app-shfl1}]
__attribute__((intel_reqd_sub_group_size(32)))
__kernel void shfl1(__global int *a, __global int *b, __global int *c, int N)
{
    int id  = get_global_id(0);
    int lix = get_sub_group_local_id();

    int sum = 0;
    for (int i = 0; i < N; i++) {
        for (int j = 0; j < N; j++)
            sum += intel_sub_group_shuffle(a[id * N + i], j);
    }
    c[lix] = sum;
}
\end{lstlisting}

\begin{lstlisting}[language=C, caption={[Source of shfl2]OpenCL source of \texttt{shfl2}.}, label={lst:app-shfl2}]
__attribute__((intel_reqd_sub_group_size(32)))
__kernel void shfl2(__global int *a, __global int *b, __global int *c, int N)
{
    int id  = get_global_id(0);
    int lix = get_sub_group_local_id();

    int sum = 0;
    for (int i = 0; i < N; i++) {
        sum += intel_sub_group_shuffle(a[id * N + i], 31);
    }
    c[lix] = sum;
}
\end{lstlisting}

\begin{lstlisting}[language=C, caption={[Source of shfl3]OpenCL source of \texttt{shfl3}.}, label={lst:app-shfl3}]
__attribute__((intel_reqd_sub_group_size(32)))
__kernel void shfl3(__global int *a, __global int *b, __global int *c, int N)
{
    int id = get_global_id(0);

    int sum = 0;
    for (int i = 0; i < N; i++) {
        int lix = (id + 1) % 32;
        sum += intel_sub_group_shuffle(a[id * N + i], lix)
             * intel_sub_group_shuffle(b[id * N + i], lix);
    }
    c[id] = sum;
}
\end{lstlisting}

\begin{lstlisting}[language=C, caption={[Source of shfl\_dn1]OpenCL source of \texttt{shfl\_dn1}.}, label={lst:app-shfl-dn1}]
__attribute__((intel_reqd_sub_group_size(32)))
__kernel void shfl_dn1(__global int *a, __global int *b, __global int *c, int N)
{
    int id  = get_global_id(0);
    int lix = get_sub_group_local_id();

    int red = a[id * N] * b[id * N];

    for (int j = 16; j > 0; j = j / 2)
        red += intel_sub_group_shuffle_down(red, red, j);
    c[lix] = red;
}
\end{lstlisting}

\begin{lstlisting}[language=C, caption={[Source of shfl\_dn2]OpenCL source of \texttt{shfl\_dn2}.}, label={lst:app-shfl-dn2}]
__attribute__((intel_reqd_sub_group_size(32)))
__kernel void shfl_dn2(__global int *a, __global int *b, __global int *c, int N)
{
    int id  = get_global_id(0);
    int lix = get_sub_group_local_id();

    for (int i = 0; i < N; i++) {
        int red = a[id * N + i] * b[id * N + i];

        for (int j = 16; j > 0; j = j / 2)
            red += intel_sub_group_shuffle_down(red, 0, j);
        if (lix == 0)
            c[i] = red;
    }
}
\end{lstlisting}

\begin{lstlisting}[language=C, caption={[Source of shfl\_dn3]OpenCL source of \texttt{shfl\_dn3}.}, label={lst:app-shfl-dn3}]
__attribute__((intel_reqd_sub_group_size(32)))
__kernel void shfl_dn3(__global int *a, __global int *b, __global int *c, int N)
{
    int id  = get_global_id(0);
    int lix = get_sub_group_local_id();

    for (int i = 0; i < N; i++) {
        int red = a[id * N + i] * b[id * N + i];

        for (int j = N / 2; j > 0; j = j / 2)
            red += intel_sub_group_shuffle_down(red, 0, j);
        if (lix == 0)
            c[i] = red;
    }
}
\end{lstlisting}

\begin{lstlisting}[language=C, caption={[Source of shfl\_up1]OpenCL source of \texttt{shfl\_up1}.}, label={lst:app-shfl-up1}]
__attribute__((intel_reqd_sub_group_size(32)))
__kernel void shfl_up1(__global int *a, __global int *b, __global int *c, int N)
{
    int id  = get_global_id(0);
    int lix = get_sub_group_local_id();

    int red = a[id * N] * b[id * N];

    for (int j = 16; j > 0; j = j / 2)
        red += intel_sub_group_shuffle_up(red, red, j);
    c[lix] = red;
}
\end{lstlisting}

\begin{lstlisting}[language=C, caption={[Source of shfl\_up2]OpenCL source of \texttt{shfl\_up2}.}, label={lst:app-shfl-up2}]
__attribute__((intel_reqd_sub_group_size(32)))
__kernel void shfl_up2(__global int *a, __global int *b, __global int *c, int N)
{
    int id  = get_global_id(0);
    int lix = get_sub_group_local_id();

    int red = a[id * N] * b[id * N];

    for (int j = N / 2; j > 0; j = j / 2)
        red += intel_sub_group_shuffle_up(red, red, j);
    c[lix] = red;
}
\end{lstlisting}

\begin{lstlisting}[language=C, caption={[Source of shfl\_up3]OpenCL source of \texttt{shfl\_up3}.}, label={lst:app-shfl-up3}]
__attribute__((intel_reqd_sub_group_size(32)))
__kernel void shfl_up3(__global int *a, __global int *b, __global int *c, int N)
{
    int id  = get_global_id(0);
    int lix = get_sub_group_local_id();

    for (int i = 0; i < N; i++) {
        int red = a[id * N + i] * b[id * N + i];

        for (int j = 16; j > 0; j = j / 2)
            red += intel_sub_group_shuffle_up(0, red, j);
        if (lix == 31)
            c[i] = red;
    }
}
\end{lstlisting}

\begin{lstlisting}[language=C, caption={[Source of shfl\_xor1]OpenCL source of \texttt{shfl\_xor1}.}, label={lst:app-shfl-xor1}]
__attribute__((intel_reqd_sub_group_size(32)))
__kernel void shfl_xor1(__global int *a, __global int *b, __global int *c, int N)
{
    int id  = get_global_id(0);
    int lix = get_sub_group_local_id();

    int sum = 0;
    for (int i = 0; i < N; i++) {
        sum += intel_sub_group_shuffle_xor(a[id * N + i], 0x1F)
             * intel_sub_group_shuffle_xor(b[id * N + i], 0x1F);
    }
    c[lix] = sum;
}
\end{lstlisting}

\begin{lstlisting}[language=C, caption={[Source of shfl\_xor2]OpenCL source of \texttt{shfl\_xor2}.}, label={lst:app-shfl-xor2}]
__attribute__((intel_reqd_sub_group_size(32)))
__kernel void shfl_xor2(__global int *a, __global int *b, __global int *c, int N)
{
    int id  = get_global_id(0);
    int lix = get_sub_group_local_id();

    if (lix < 16) {
        int sum = 0;
        for (int i = 0; i < N; i++) {
            sum += intel_sub_group_shuffle_xor(a[id * N + i], 0xF)
                 * intel_sub_group_shuffle_xor(b[id * N + i], 0xF);
        }
        c[lix] = sum;
    }
}
\end{lstlisting}

\begin{lstlisting}[language=C, caption={[Source of prefix\_sum]OpenCL source of \texttt{prefix\_sum}.}, label={lst:app-prefix-sum}]
__kernel void prefix_sum(__global const int *a, __global int *b)
{
    uint gid = get_global_id(0);
    int x = a[gid];

    b[gid] = sub_group_scan_inclusive_add(x);
}
\end{lstlisting}

\begin{lstlisting}[language=C, caption={[Source of ired\_sum]OpenCL source of \texttt{ired\_sum}.}, label={lst:app-ired-sum}]
__kernel void ired_sum(__global int *a, __global int *c)
{
    int id = get_global_id(0);
    c[id] = sub_group_reduce_add(a[id]);
}
\end{lstlisting}

\begin{lstlisting}[language=C, caption={[Source of reduce\_if\_else]OpenCL source of \texttt{reduce\_if\_else}.}, label={lst:app-reduce-if-else}]
__kernel void reduce_if_else(__global const int *a, __global int *out)
{
    int gid = get_global_id(0);
    int lid = get_sub_group_local_id();
    if (lid < 16)
        out[gid] = sub_group_reduce_add(a[gid]);
    else
        out[gid] = -1;
}
\end{lstlisting}

\section*{Memory}

\begin{lstlisting}[language=C, caption={[Source of exchange]OpenCL source of \texttt{exchange}.}, label={lst:app-exchange}]
__kernel void exchange(__global const int *a, __global int *out)
{
    __local int buf[32];
    int gid = get_global_id(0);
    int lid = get_sub_group_local_id();
    buf[lid] = a[gid];
    out[gid] = buf[31 - lid];
}
\end{lstlisting}